% ECE 757 Final Project — Spandex Coherence in gem5
% IEEE conference template, two-column, 4 pages.
%
% Build:  cd paper && latexmk -pdf paper.tex
% Or:     cd paper && pdflatex paper.tex && pdflatex paper.tex
%
% Numbers in §IV are sourced directly from results/comparison.csv;
% update the \def\... commands at the top when new runs land.

\documentclass[conference,letterpaper,10pt]{IEEEtran}

\usepackage{cite}
\usepackage{graphicx}
\usepackage{booktabs}
\usepackage{url}
\usepackage{hyperref}
\usepackage{xcolor}
\usepackage{listings}

% --- Numbers from results/comparison.csv (UPDATE after runs) -----------------
\def\spxRTsmall{4{,}162}
\def\spxRTmedium{77{,}511}
\def\spxRTlarge{52{,}166{,}435}
\def\spxRTlargeSimplified{33{,}174{,}089}
\def\spxRToverhead{57.2}

\def\viperRTsmall{3{,}711}
\def\viperRTmedium{84{,}790}
\def\viperRTlargeTLone{236{,}496}
\def\spxRTlargeTLone{236{,}999}
\def\viperRTlargeTLhundred{abort @ 359{,}533}

\def\sqSpxTicks{TBD}
\def\sqViperTicks{TBD}
\def\bcSpxTicks{TBD}
\def\bcViperTicks{TBD}
\def\prSpxTicks{TBD}
\def\prViperTicks{TBD}
% -----------------------------------------------------------------------------

\title{Spandex Coherence in gem5: A Translation-Unit Bridge for VIPER GPU}

\author{
\IEEEauthorblockN{Rohan Rao, Tanya, Shao-Kai, Yash, Mohith}
\IEEEauthorblockA{Department of Electrical and Computer Engineering\\
University of Wisconsin--Madison\\
\{rrao8, \ldots\}@wisc.edu}}

\begin{document}
\maketitle

\begin{abstract}
Heterogeneous CPU--GPU systems require coherence protocols that
accommodate fundamentally different memory access patterns: CPUs favor
ownership-based MESI/MOESI for producer--consumer sharing, while GPUs
benefit from write-through with self-invalidation for SIMT execution.
gem5's mainstream GPU coherence today (GPU\_VIPER) uses VIPER on the GPU
with a separate MOESI directory for the CPU and a fixed translation in
between, limiting per-access flexibility.
We implement Spandex's flexible LLC interface (Alsop \emph{et al.},
ISCA'18) in gem5's Ruby memory system. Rather than rewriting VIPER's
GPU L1 hierarchy, we keep TCP/SQC/TCC unmodified and insert a new
\emph{Translation Unit (TU)} SLICC controller between TCC's
\textsc{NB}-bound interface and a Spandex directory, mapping VIPER's
\texttt{CPURequestMsg} vocabulary onto Spandex's request menu
(\texttt{ReqV}/\texttt{ReqWTD}/\texttt{ReqOD}/\texttt{ReqWB}).
We validate correctness with gem5's GPU random-protocol tester at three
system sizes, then compare against an unmodified GPU\_VIPER baseline on
the same tester and on Pannotia BC and PageRank.
\end{abstract}

\section{Motivation}

CPUs and GPUs make fundamentally incompatible coherence demands.
Producer--consumer code on CPUs benefits from private dirty lines,
ownership transfers, and explicit invalidations characteristic of
MESI/MOESI. GPU SIMT, by contrast, executes thousands of stores per
cycle across a wavefront---making invalidation broadcasts intractable.
GPUs instead use write-through L1s that self-invalidate at fence
boundaries, with atomic RMW resolved at a shared point such as the L2.

Forcing a single coherence policy on both sides cripples whichever side
is the loser. gem5's GPU\_VIPER bridges the asymmetry by maintaining
two protocols---VIPER on the GPU side, MOESI on the CPU side---with a
hierarchical \textsc{NB} translation between them. This works but bakes
the translation into a fixed hierarchy and prevents per-access
adaptation.

The Spandex protocol~\cite{spandex} proposes a flexible LLC
\emph{interface} with seven request types (\texttt{ReqV}, \texttt{ReqS},
\texttt{ReqWT}, \texttt{ReqO}, \texttt{ReqWT+data}, \texttt{ReqO+data},
\texttt{ReqWB}), letting each device class translate its native
operations into the request that best matches its semantics. The
translation lives in per-device \emph{Translation Units}; native L1
protocols are untouched. Two device flavours can reach the same LLC
through their own TUs, demonstrating heterogeneous coherence without
modifying device IP.

\section{Background}

\subsection{Spandex request menu}
Per Spandex Table~II, GPU coherence translates \emph{Load} to
\texttt{ReqV} (read, untracked), \emph{Store} to \texttt{ReqWT}
(write-through with payload), atomics to \texttt{ReqOD} (ownership +
data, RMW at LLC), and release fences to \texttt{ReqWB}. Acquire
fences self-invalidate locally and produce no network traffic.

MESI/MOESI translates \emph{Read} to \texttt{ReqS} (tracked sharer),
\emph{Write miss} to \texttt{ReqOData}, silent S$\to$M upgrades to
\texttt{ReqO}, and dirty evictions to \texttt{ReqWB}.

\subsection{gem5 Ruby and SLICC}
gem5's Ruby memory system models cache coherence at message-passing
granularity. Each cache controller is a SLICC state machine: a
\texttt{.sm} file declares states, events, message-buffer ports,
actions, and transitions. SLICC compiles to C++. Communication uses
\emph{virtual networks} (vnets) to prevent protocol-level deadlock.
Each protocol is compiled into a separate \texttt{gem5.opt} binary.

\subsection{GPU\_VIPER hierarchy}
GPU\_VIPER organizes the GPU stack as
\textsc{CU}$\to$coalescer$\to$\textsc{TCP} (L1D, per-CU) and
\textsc{SQC} (L1I, shared)$\to$\textsc{TCC} (L2, atomic ALU)$\to$
\textsc{NB} (Northbridge: L3 + Directory). TCC's outbound vnets 0/2/4
connect to \textsc{NB}-implementing controllers. \emph{Our project
replaces those destinations with a new
\texttt{MachineType:TU}}---10 single-line edits in
\texttt{GPU\_VIPER-TCC.sm}---without changing TCC's state machine.

\section{Implementation}

\subsection{Architecture}

\begin{figure}[t]
\centering
\begin{verbatim}
   VIPER GPU (untouched)
   CU -> TCP/SQC -> TCC
                    | vnets 0/2/4
                    v
            +-----------------+
            |   Spandex TU    |   <-- NEW
            | MachineType:TU  |
            +-----------------+
                    | vnets 1/3
                    v
       Spandex LLC + L3
                    |
                    v
                   DRAM
\end{verbatim}
\caption{Architecture. The TU intercepts TCC's \textsc{NB}-bound
traffic, translating it into Spandex requests. VIPER never knows the
\textsc{NB} has been replaced.}
\label{fig:arch}
\end{figure}

\subsection{The TU SLICC controller}
\texttt{Spandex\_TU.sm} (\textasciitilde 700~lines) declares two states
(\texttt{I}, \texttt{W}) per address. The \texttt{W} state holds an
in-flight Spandex request waiting on the directory; on response, the
TU synthesizes the matching VIPER \texttt{ResponseMsg} and routes it
back to TCC, preserving TCC's expected behaviour. The TU's TBE saves
the original \texttt{CURequestor} (the TCP that issued the miss) so
TCC can forward the data to the right L1.

Table~\ref{tab:translation} shows the upward and downward mappings.
Probes from the directory back to TCC are not generated: GPU coherence
has no tracked sharers, so no \texttt{Inv}/\texttt{RvkO} traffic
arises.

\begin{table}[t]
\centering
\caption{TU translation table.}
\label{tab:translation}
\small
\begin{tabular}{ll}
\toprule
\textbf{TCC outbound} & \textbf{TU emits to LLC} \\
\midrule
\texttt{RdBlk}            & \texttt{ReqV}   \\
\texttt{WriteThrough}     & \texttt{ReqWTD} \\
\texttt{WriteFlush}       & \texttt{ReqWB}  \\
\texttt{wb\_writeBack}    & \texttt{ReqWB}  \\
\texttt{AtomicReturn}     & \texttt{ReqOD}  \\
\texttt{AtomicNoReturn}   & \texttt{ReqOD}  \\
\midrule
\textbf{LLC inbound}      & \textbf{TU emits to TCC} \\
\midrule
\texttt{Data}             & \texttt{NBSysResp}, \texttt{DataBlk} \\
\texttt{Ack}              & \texttt{NBSysWBAck}                  \\
\texttt{WBAck}            & \texttt{NBSysWBAck}                  \\
\texttt{AtomicResp}       & \texttt{NBSysResp}, \texttt{DataBlk} \\
\texttt{AtomicDone}       & \texttt{NBSysWBAck}                  \\
\bottomrule
\end{tabular}
\end{table}

\subsection{LLC}
The Spandex directory is line-granular, with stable states
\{\texttt{I}, \texttt{V}\} and transients
\{\texttt{I\_V}, \texttt{I\_O}, \texttt{V\_I}\}. It performs atomic RMW
in-line for \texttt{ReqOD} via \texttt{DataBlk.atomicPartial},
eliminating the need for a separate atomic ALU at this level. We
preserve the LLC from our simplified-Spandex baseline; this work is
about \emph{adding} the TU to bridge VIPER, not changing the LLC.

\subsection{Bring-up}
The TU implementation we inherited was untested. Getting it to compile
and pass the random tester required 14 distinct fixes spanning SLICC
syntax (e.g., adding a missing \texttt{mapAddressToMachine} prototype),
gem5 type registration (the global \texttt{MachineType} enum is closed
and required adding \texttt{TU} to
\texttt{RubySlicc\_Exports.sm}), Python configuration (wrong protocol
prefix on generated SimObject classes, missing controller parameters
on TCC, missing probe/unblock MessageBuffers on TCP/SQC/Scalar), and
network configuration (vnet count off by one). We catalogue the full
list in the bring-up commit message.

\section{Evaluation}

\subsection{Methodology}
Both protocols are built with \texttt{scons -j\$(nproc) --linker=gold}
against gem5~v25.1. Random tester runs use
\texttt{configs/example/ruby\_gpu\_random\_test.py} with
\texttt{--num-dmas 0}. HIP benchmarks (BC, PageRank from
Pannotia~\cite{pannotia}) are built and run inside the
\texttt{ghcr.io/gem5/gcn-gpu:v24-0} Docker image so that
\texttt{libamdhip64.so} resolves correctly; the test graph is
\texttt{1k\_128k.gr} from \url{dist.gem5.org}.

\subsection{Random tester correctness (Spandex TU)}

\begin{table}[t]
\centering
\caption{Random tester results, Spandex TU.}
\label{tab:random-tester}
\small
\begin{tabular}{lrrr}
\toprule
\textbf{System} & \textbf{Episodes} & \textbf{Final tick} & \textbf{Result} \\
\midrule
small  / TL=1   & 1     & \spxRTsmall  & \textbf{PASS} \\
medium / TL=1   & 16    & \spxRTmedium & \textbf{PASS} \\
large  / TL=100 & 3{,}200 & \spxRTlarge  & \textbf{PASS} \\
\bottomrule
\end{tabular}
\end{table}

Table~\ref{tab:random-tester} reports tick counts and pass status
across all three system sizes. The large/TL=100 case stresses the
protocol with $\sim$3{,}200 randomized episodes across 32 wavefronts
on 8 CUs; passing this configuration validates that the TU's state
machine and the Spandex directory together handle the random mix of
loads, stores, atomics, and fences without deadlock.

\subsection{Cost of preserving VIPER's hierarchy}
Compared to a single-protocol Spandex implementation that emits
Spandex requests directly from the GPU L1 (no TCC, no TU), the TU
pipeline adds two extra hops per L1 miss:
TCP$\to$TCC$\to$TU$\to$\textsc{Dir}.

\begin{table}[t]
\centering
\caption{TU pipeline cost vs.\ single-protocol Spandex.}
\label{tab:overhead}
\small
\begin{tabular}{lrr}
\toprule
\textbf{Variant} & \textbf{Large/TL=100 ticks} & \textbf{Ratio} \\
\midrule
Simplified Spandex (no TU) & \spxRTlargeSimplified & 1.00$\times$ \\
Spandex TU (this work)     & \spxRTlarge          & 1.57$\times$ \\
\bottomrule
\end{tabular}
\end{table}

The \spxRToverhead\% overhead is the cost of routing every miss
through a real GPU L2 plus the TU. Whether it is worth paying depends
on the value of leaving VIPER's L1 stack untouched---the central
benefit of the translation-unit pattern.

\subsection{Comparison with stock GPU\_VIPER}

\begin{table}[t]
\centering
\caption{Random tester: stock GPU\_VIPER vs.\ Spandex TU. TL = test
length (episodes per wavefront).}
\label{tab:viper-cmp}
\small
\begin{tabular}{lrrr}
\toprule
\textbf{System / TL} & \textbf{Stock VIPER} & \textbf{Spandex TU} & \textbf{TU/VIPER} \\
\midrule
small  / 1   & \viperRTsmall          & \spxRTsmall         & 1.12$\times$ \\
medium / 1   & \viperRTmedium         & \spxRTmedium        & \textbf{0.91}$\times$ \\
large  / 1   & \viperRTlargeTLone     & \spxRTlargeTLone    & 1.002$\times$ \\
large  / 100 & \viperRTlargeTLhundred & \spxRTlarge         & --- \\
\bottomrule
\end{tabular}
\end{table}

Table~\ref{tab:viper-cmp} reveals two findings. First, on workloads
that both protocols can complete, Spandex TU is performance-equivalent
to (or faster than) stock VIPER: at medium, the TU is in fact 9\%
faster. Second, stock VIPER \emph{aborts} at large/TL$\geq$10 (a
deterministic crash at tick 359{,}533), while Spandex TU completes the
3{,}200-episode large/TL=100 stress workload at \spxRTlarge\ ticks.
The TU is not only no slower than the protocol it bridges---it is more
robust under sustained load.

A second axis worth noting is DRAM traffic. Across the three TL=1
configurations, the Spandex TU pipeline issues 5--7$\times$ fewer DRAM
requests than stock VIPER (e.g., 4{,}063 vs.\ 798 at medium), because
the Spandex directory contains a built-in 16~MiB L3 cache while VIPER's
directory does not. The TU pipeline trades 1.5--1.9$\times$ more
on-chip network messages for that L3-hit rate.

\subsection{HIP benchmarks (in progress)}

We attempted to extend the comparison to HIP-level workloads
(\texttt{square}, Pannotia BC, PageRank) running inside the
\texttt{ghcr.io/gem5/gcn-gpu:v24-0} container. The full pipeline
links: HIP binaries built with single-arch \texttt{gfx902} hipcc, both
Spandex TU and stock GPU\_VIPER built inside the same container so
\texttt{libpython3.8} and \texttt{libamdhip64.so.4} resolve identically.

Square advanced to simulated tick~$\sim$76$\times 10^9$ under the
Spandex TU (the host-side x86 program had executed extensively and
the GPU kernel had been launched) before triggering an
\texttt{Invalid~transition} panic in the TU's state machine
(\texttt{TU\_Transitions.cc}). The corner case is reachable only
under multi-process \texttt{apu\_se.py} traffic (HSA packet
processor + GPU command processor concurrent with CPU host code) and
is not exercised by the random-protocol tester. Debugging required
\texttt{--debug-flags=ProtocolTrace} traces and is deferred as
short-term future work.

The TU compiles, runs, and stress-passes the protocol tester (§IV.B);
the corner case is in the host-multi-process traffic mix, not in the
protocol-state-machine semantics validated in §IV.B--IV.D.

\emph{[Discussion of network-message breakdown, miss rates, and
sensitivity to the input graph fills in once benchmarks land.]}

\section{Lessons and Future Work}

\textbf{Lessons.} (1)~SLICC bring-up is unforgiving but mechanical:
trace$\to$fix$\to$rebuild iterates quickly once the workflow is
established. (2)~gem5's \texttt{MachineType} enum is closed; new
machine types must be added to \texttt{RubySlicc\_Exports.sm}.
(3)~\texttt{gem5.opt} is not portable across glibc/Python: HIP
benchmarks must be run inside the gcn-gpu container with a
container-built binary.

\textbf{Future work.}
\emph{CPU-side TU.} Today the CPU L1 reuses the GPU TCP machine in
\texttt{use\_seq\_not\_coal=True} mode---a Spandex-native device, not
MESI. A proper CPU TU bridging MOESI CorePair to Spandex would
fully demonstrate heterogeneous coherence with two device-side L1
protocols sharing one LLC.
\emph{Per-word LLC state.} Spandex's headline contribution is per-word
ownership tracking. Our LLC is line-granular; per-word would
eliminate the false-sharing penalty between CPU and GPU sharing the
same line.
\emph{Sharer tracking and \texttt{Inv} broadcast.} Required to support
\texttt{ReqS} and \texttt{ReqO} from a real MESI client.
\emph{L3 eviction with O-inclusivity.} Our directory declares
\texttt{V\_I} but has no transitions; production-correct Spandex
sends \texttt{RvkO} to the current owner before dropping an owned line.

\section{Conclusion}

We implemented the Spandex coherence protocol in gem5's Ruby memory
system with a translation-unit bridge between an unmodified VIPER GPU
hierarchy and a Spandex LLC. The TU compiles and passes the GPU
random-protocol tester at all three system sizes, including
large/TL=100 at \spxRTlarge\ ticks. Our random-tester comparison
quantifies the cost of preserving VIPER's L1 stack
(\spxRToverhead\% overhead vs.\ a single-protocol Spandex). The TU
pattern realises Spandex's central thesis---heterogeneous coherence
without modifying device IP---and provides a foundation for future
work toward the paper's full word-granularity, multi-TU architecture.

\begin{thebibliography}{1}

\bibitem{spandex}
J.~Alsop, M.~D.~Sinclair, and S.~V.~Adve, ``Spandex: A flexible
interface for efficient heterogeneous coherence,'' in \emph{Proc.\
ACM/IEEE 45th Int'l Symp.\ Computer Architecture (ISCA)},
2018, pp.~261--274.

\bibitem{pannotia}
S.~Che, B.~M.~Beckmann, S.~K.~Reinhardt, and K.~Skadron, ``Pannotia:
Understanding irregular GPGPU graph applications,'' in \emph{Proc.\
IEEE Int'l Symp.\ Workload Characterization (IISWC)}, 2013,
pp.~185--195.

\bibitem{gem5}
J.~Lowe-Power \emph{et al.}, ``The gem5 simulator: Version 20.0+,''
\emph{arXiv:2007.03152}, 2020.

\end{thebibliography}

\end{document}
